% Generated draft: requires mathematical and source review.
\documentclass[11pt]{article}
\usepackage[margin=25mm]{geometry}
\usepackage{fontspec}
\setmainfont{Latin Modern Roman}
\usepackage{amsmath,amssymb,mathtools,unicode-math}
\setmathfont{Latin Modern Math}
\usepackage{enumitem,needspace,xurl,hyperref}
\setlength{\parindent}{0pt}
\setlength{\parskip}{5pt}
\setlength{\emergencystretch}{3em}
\newcommand{\N}{\mathbb{N}}
\newcommand{\Z}{\mathbb{Z}}
\newcommand{\Q}{\mathbb{Q}}
\newcommand{\R}{\mathbb{R}}
\newcommand{\C}{\mathbb{C}}
\newcommand{\F}{\mathbb{F}}
\newcommand{\A}{\mathbb{A}}
\newcommand{\PP}{\mathbb{P}}
\newcommand{\E}{\mathbb{E}}
\newcommand{\eps}{\varepsilon}
\newcommand{\dd}{\,\mathrm{d}}
\begin{document}
\section*{UnsolvedMath ID: 20000560}
\textbf{Source problem code:} catalogue entry
\subsection{Definitions and mathematical statement}
The AIM fragment asks whether a Kadison--Singer-related partition can improve the recorded $n^{5/6}$ Banach--Mazur distance bound from a convex body to the cube. Its phrase ``subspaces like this'' refers to missing context. For a normed-space formulation, take $X$ a real $n$-dimensional space with an origin-symmetric convex unit ball and
\[
 d_{\mathrm{BM}}(X,\ell_\infty^n)=\inf_{T:X\cong\ell_\infty^n}\|T\|\,\|T^{-1}\|.
\]
The original does not state symmetry, normalization or the partition's precise structure. Its best-bound phrase is a historical source claim without a supplied date or reference.

The report's gluing theorem requires a genuine linear direct sum $X=\bigoplus_{j=1}^rX_j$, with dimensions $d_j$ summing to $n$, and two-sided aggregation bounds, for $1\leq q\leq\infty$,
\[
 a\left(\sum_j\|x_j\|^q\right)^{1/q}
 \leq\left\|\sum_jx_j\right\|_X
 \leq b\left(\sum_j\|x_j\|^q\right)^{1/q},\qquad0<a\leq b,
\]
using maxima for $q=\infty$. If $d_{\mathrm{BM}}(X_j,\ell_\infty^{d_j})\leq D$ for every $j$, it claims
\[
 d_{\mathrm{BM}}(X,\ell_\infty^n)\leq(b/a)\,r^{1/q}D.
\]
For balanced blocks and quadratic aggregation, combining local bounds of size $C d_j^{5/6}$ yields
\[
 d_{\mathrm{BM}}(X,\ell_\infty^n)\leq C(b/a)n^{5/6}r^{-1/3}.
\]
Along a sequence this particular upper bound is $o(n^{5/6})$ exactly when $b/a=o(r^{1/3})$. This is a criterion on the reduction's estimate, not a necessary condition for every actual improvement by other methods.

The report also gives a determinant-volume obstruction. For $X=(\bigoplus_{j=1}^r\ell_\infty^d)_{\ell_1}$ with $n=rd$, every coordinate block section is exactly a cube, yet
\[
 d_{\mathrm{BM}}(X,\ell_\infty^{rd})\geq
 \left(\frac{(rd)!}{r^{rd/2}(d!)^r}\right)^{1/(rd)}
 \geq\frac{\sqrt r}{e}.
\]
This is a lower bound, not a formula for the exact global distance. Distortion-one local sections therefore do not eliminate a growing global aggregation cost.

Finally, the report notes that if each frame block satisfies two-sided positive ambient bounds $aI\leq\sum_{i\in J}v_iv_i^T\leq bI$, then each block spans the entire ambient space. Such a frame partition is not automatically a complementary subspace decomposition of the kind required above. No construction of the needed uniformly controlled direct sum from Kadison--Singer, and no improvement of the universal exponent, is claimed.
\subsection{Short English statement}
Ask whether Kadison--Singer partitions can improve cube-distance bounds. The report identifies the direct-sum and norm-aggregation control that a gluing argument would need, gives an explicit quadratic threshold, and shows why cube-like sections alone are insufficient. It does not supply a new universal exponent or turn frame blocks automatically into complementary subspaces.
\subsection{Sources}
\begin{itemize}
\item Local registry: \url{data/unsolvedmath/problems.json}, permanent ID $20000560$, source code \texttt{catalogue entry}; original fragment and supplied research summary.
\item AIM workshop ``Beyond Kadison-Singer: paving and consequences'', source item b19, Google document: \url{https://docs.google.com/document/d/1XRw0WyG5PshqboCp99tovCRD1GCK0jXZiNLH_nw618c/edit}. Canonical location \texttt{aim-analysis-notes.json}, \texttt{notes[91]}; original remarks and literature field empty. No document was visited here.
\item No full authoritative citation or date is supplied for the historical $n^{5/6}$ benchmark, gluing theorem or factorial obstruction.
\item Supplied machine-generated research attempt 1: reduction, problem open at the run and currently partially solved; rigor metadata ``proved'', novelty confidence low. This non-peer-reviewed synthesis requires independent verification before citation.
\end{itemize}
\end{document}
